\documentclass{article}
\usepackage[T1]{fontenc}
\usepackage{lmodern}
\usepackage{graphicx}
\usepackage{amsmath,amssymb}
\usepackage[a4paper,margin=2cm]{geometry}
\usepackage[absolute,overlay]{textpos}
\setlength{\TPHorizModule}{1mm}
\setlength{\TPVertModule}{1mm}
\usepackage[indonesian]{babel}
\usepackage{hyperref}
\setlength{\emergencystretch}{2em}
% unit: Halaman 5

\begin{document}

% === Halaman 5 (python/native) ===

Contoh: Misalkan plainteks adalah 110111010001000101000101, dibagi menjadi 

tiga buah blok masing-masing berukuran 8 bit:


 11011101     00010001      01000101

Kunci yang digunakan adalah K = 01010110 

Enkripsi blok pertama, 11011101,  dengan fungsi E sebagai berikut:


 1.  P1  K = 11011101   01010110 = 10001011


 2.  10001011 << 1 = 00010111

Jadi,  C1 = 00010111. Cara yang sama dilakukan untuk P2 dan P3.

Dekripsi blok cipherteks pertama, 00010111, dengan fungsi D sebagai berikut:


1.  00010111 >> 1 = 10001011 


2.  10001011  K = 10001011  01010110 = 11011101 = P1

5

P1 

          P2  

        P3

\end{document}
